\subsection{从无穷小作用律到作用律的过渡}

命题 11. 设 \(\mathrm{G}_1\) 和 \(\mathrm{G}_2\) 是李群芽，\(\mathbf{X}_1\) 和 \(\mathbf{X}_2\) 是 \(\mathrm{C}^r\) 类流形（\(r \geqslant 2\)）。对于 \(i = 1, 2\)，令 \(\psi_i\) 为 \(\mathrm{C}^r\) 类 \(\mathrm{G}_i\) 在 \(\mathbf{X}_i\) 上的左作用律片段，\(\mathrm{D}_i\) 为相应的无穷小作用律。令 \(\mu: \mathrm{G}_1 \to \mathrm{G}_2\) 为态射，\(\phi: \mathrm{X}_1 \to \mathrm{X}_2\) 为 \(\mathrm{C}^r\) 类映射。假设对所有 \(a \in \mathrm{L}(\mathrm{G})\)，向量场 \((\mathrm{D}_{1})_a\) 和 \((\mathrm{D}_{2})_{\mathrm{L}(\mu)a}\) 关于 \(\phi\) 相关（Differentiable and Analytic Manifolds, R, 8.2.6）。则存在一个邻域 \(\Omega\)，它包含 \(\{e\} \times \mathbf{X}_1\) 且位于 \(\mathrm{G}_1 \times \mathbf{X}_1\) 中，使得 \(\phi(\psi_1(g,x)) = \psi_2(\mu(g),\phi(x))\) 对所有 \((g,x) \in \Omega\) 成立。

令 \(p_1: \mathrm{G}_1 \times \mathrm{X}_2 \to \mathrm{G}_1\)、\(p_2: \mathrm{G}_1 \times \mathrm{X}_2 \to \mathrm{X}_2\) 为规范投影。对所有 \((g_1, x_2) \in \mathrm{G}_1 \times \mathrm{X}_2\)，令 \(f_{g_1, x_2}\) 为映射 \(g_1 a \mapsto (\mathrm{D}_2)_{\mathrm{L}(\mu)a}(x_2)\)，它从 \(\mathrm{T}_{g_1}(\mathrm{G}_1) = g_1\mathrm{L}(\mathrm{G}_1)\) 到 \(\mathrm{T}_{x_2}(\mathrm{X}_2)\)。这些 \(f_{g_1, x_2}\) 定义了从 \(p_1^* \mathrm{T}(\mathrm{G}_1)\) 到 \(p_2^* \mathrm{T}(\mathbf{X}_2)\) 的态射。

令 \(x_0 \in X_1\)。存在 e 的一个开邻域 G（位于 \(G_1\) 中）以及 \(x_0\) 在 \(X_1\) 中的一个开邻域 X，使得 \(\psi_1(g, x)\) 和 \(\psi_2(\mu(g), \phi(x))\) 对 \((g, x) \in G \times X\) 有定义。对于 \((g, x) \in G \times X\)，记

\[
\alpha (g, x) = \phi (\psi_ {1} (g, x)) \in \mathrm{X} _ {2}, \quad \beta (g, x) = \psi_ {2} (\mu (g), \phi (x)) \in \mathrm{X} _ {2}.
\]

若 \(\mathbf{G}\) 和 \(\mathbf{X}\) 足够小，则对所有 \((a,g,x)\in \mathrm{L}(\mathrm{G}_1)\times \mathrm{G}\times \mathrm{X},\)

\[
\begin{array}{r l} (\mathrm{T} \alpha) (a g, 0 _ {x}) & = (\mathrm{T} \phi) ((\mathrm{D} _ {1}) _ {a} (\psi_ {1} (g, x))) \\ & = (\mathrm{D} _ {2}) _ {L (\mu) a} (\phi (\psi_ {1} (g, x))) \\ & = (\mathrm{D} _ {2}) _ {L (\mu) a} \alpha (g, x), \end{array}
\]

\[
\begin{array}{r l} (\mathrm{T} \beta) (a g, 0 _ {x}) = & (\mathrm{T} \psi_ {2}) (\mathrm{L} (\mu) a. \mu (g), \phi (x)) \\ = & (\mathrm{D} _ {2}) _ {\mathrm{L} (\mu) a} (\psi_ {2} (\mu (g), \phi (x))) \\ = & (\mathrm{D} _ {2}) _ {\mathrm{L} (\mu) a} \beta (g, x). \end{array}
\]

因此对 \(x \in \mathbf{X}\)，映射 \(g \mapsto \alpha(g, x)\) 和 \(g \mapsto \beta(g, x)\) 都是 \(f\) 的积分；由于

\[
\beta (e, x) = \phi (x) = \alpha (e, x)
\]

对所有 \(x \in \mathbf{X}\)，由 Differentiable and Analytic Manifolds, R, 9.3.7，可知 \(\alpha\) 和 \(\beta\) 在 \((e, x_0)\) 的某个邻域上重合。因此命题得证。

推论。设 G 是李群芽，X 是 \(\mathbf{C}^{r}\) 类流形。考虑 G 在 X 上的两个 \(\mathbf{C}^{r}\) 类左作用律片段。假设对所有 \(a \in L(G)\)，相应的 X 上的向量场 \(D_{a}\) 对两个作用律片段都相同。则这两个作用律片段在 \(\{e\} \times X\) 的某个邻域上重合。

定理 6. 设 G 是李群芽，X 是 \(\mathbf{C}^r\) 类流形（\(r \geqslant 2\)），\(x_0\) 是 X 中的一点。令 \(a \mapsto \mathrm{D}_a\) 为一个 \(\mathbf{C}^{r-1}\) 类的、在 X 上的 \(\mathrm{L(G)}\) 左无穷小作用律。

\begin{enumerate}
\def\labelenumi{(\roman{enumi})}
\item
  存在一个开邻域 \(\mathbf{X}'\)，它包含 \(x_0\) 且位于 \(\mathbf{X}\) 中，并且存在一个 \(\mathbf{C}^{r-1}\) 类的、由 \(\mathbf{G}\) 在 \(\mathbf{X}'\) 上定义的左作用律片段，使其相应的无穷小作用律为 \(a \mapsto \mathbf{D}_a|\mathbf{X}'\)。
\item
 设有两个 \(\mathbf{C}^{r-1}\) 类的、由 \(\mathbf{G}\) 在 \(\mathbf{X}''\)（它是 \(x_0\) 的一个开邻域）上定义的左作用律片段；若它们都以 \(a \mapsto \mathbf{D}_a|\mathbf{X}''\) 为相应的无穷小作用律，则它们在 \((e, x_0)\) 的某个邻域上重合。
\end{enumerate}

断言 (ii) 由命题 11 的推论推出。下面证明 (i)。对所有 \((g, x) \in \mathbf{G} \times \mathbf{X}\) 和 \(a \in \mathrm{L}(\mathbf{G})\)，记

\[
\mathrm{Q} _ {a} (g, x) = (a g, \mathrm{D} _ {a} (x)) \in \mathrm{T} _ {g} (\mathrm{G}) \times \mathrm{T} _ {x} (\mathrm{X}).
\]

令 \(S_{(g,x)}\) 为由 \(Q_{a}(g,x)\)（其中 \(a\in L(G)\)）组成的集合。由第 6 号的引理 5，\(S_{(g,x)}\) 是 \(T(G)\times T(X)\) 的一个向量子丛 S 的纤维。令 a, b 属于 \(L(G)\)；则

\[
\begin{array}{r l} {[ \mathrm{Q} _ {a}, \mathrm{Q} _ {b} ] (g, x) = ([ \mathrm{R} _ {a}, \mathrm{R} _ {b} ] (g), [ \mathrm{D} _ {a}, \mathrm{D} _ {b} ] (x))} \\ & = (- \mathrm{R} _ {[ a, b ]} (g), - \mathrm{D} _ {[ a, b ]} (x)) \\ & = \mathrm{Q} _ {- [ a, b ]} (g, x) \end{array}\tag{§ 3, 18.6}
\]

从而 S 可积（Differentiable and Analytic Manifolds, R, 9.3.3, (iv)）。由 Differentiable and Analytic Manifolds, R, 9.3.7，存在一个开邻域 \(\mathrm{G}_1\)，它是 G 中 \(e\) 的邻域；一个开邻域 \(\mathbf{X}_1\)，它是 X 中 \(x_0\) 的邻域；以及一个映射 \((g,x)\mapsto gx\)，它是 \(\mathrm{C}^{r - 1}\) 类，从 \(\mathrm{G}_1\times \mathrm{X}_1\) 映到 X，使得 \(ex = x\) 对所有 \(x\in \mathbf{X}_1\) 都成立，并且

\[
(a g) x = \mathrm{D} _ {a} (g x) \quad \text { for } a \in \mathrm{L} (\mathrm{G}), g \in \mathrm{G} _ {1}, x \in \mathrm{X} _ {1}.\tag{6}
\]

特别地

\[
a x = \mathrm{D} _ {a} (x).\tag{7}
\]

令 \(G_{2}\) 为 \(G_{1}\) 中 e 的一个开邻域，\(X_{2}\) 为一个开邻域，它包含 \(x_{0}\) 且位于 \(X_{1}\) 中，使得 \(gg'\) 有定义且属于 \(G_{1}\)，其中 g、\(g'\) 属于 \(G_{2}\)，并且 gx 有定义且属于 \(\mathbf{X}_1\)，其中 \((g,x)\in \mathbf{G}_2\times \mathbf{X}_2\)。考虑映射 \(\alpha_{1},\alpha_{2}\)，它们从 \(\mathrm{G}_2\times (\mathrm{G}_2\times \mathrm{X}_2)\) 到 X，定义为

\[
\alpha_ {1} (g, (h, x)) = g (h x), \quad \alpha_ {2} (g, (h, x)) = (g h) x.
\]

它们属于 \(\mathbf{C}^{r - 1}\) 类。于是

\[
\alpha_ {1} (e, (h, x)) = h x = \alpha_ {2} (e, (h, x)).
\]

另一方面

\[
\begin{array}{r l} \mathrm{T} (\alpha_ {1}) (a g, 0 _ {(h, x)}) & = (a g) (h x) \\ & = \mathrm{D} _ {a} (g (h x)) \qquad \text {by (6)} \\ & = \mathrm{D} _ {a} (\alpha_ {1} (g, (h, x))), \\ \mathrm{T} (\alpha_ {2}) (a g, 0 _ {(h, x)}) & = (a g h) x \\ & = \mathrm{D} _ {a} ((g h) x) \qquad \text {by (6)} \\ & = \mathrm{D} _ {a} (\alpha_ {2} (g, (h, x))). \end{array}
\]

由 Differentiable and Analytic Manifolds, R, 9.3.7，\(\alpha_{1}\) 和 \(\alpha_{2}\) 在 \((e, (e, x_{0}))\) 的某个邻域上重合。于是 (i) 由 § 1, no. 11 的 Proposition 23 和 (7) 得出。

推论 1. 设 G 是李群芽，X 是 \(\mathbf{C}^r\) 类旁紧流形（\(r \geqslant 2\)）。令 \(a \mapsto \mathrm{D}_a\) 为 L(G) 在 X 上的 \(\mathbf{C}^{r-1}\) 类左无穷小作用律。

\begin{enumerate}
\def\labelenumi{(\roman{enumi})}
\item
  存在一个 \(\mathbf{C}^{r-1}\) 类的、由 \(\mathbf{G}\) 在 \(\mathbf{X}\) 上定义的左作用律片段，使其相应的无穷小作用律为 \(a \mapsto \mathbf{D}_a\)。
\item
  两个 \(\mathbf{C}^{r-1}\) 类的 G 在 X 上的左作用律，若都以 \(a \mapsto \mathbf{D}_a\) 为相应的无穷小作用律，则它们在 \(\{e\} \times \mathbf{X}\) 的某个邻域上重合。
\end{enumerate}

断言 (ii) 由命题 11 的推论推出。由定理 6 (i)，存在一个开覆盖 \((\mathbf{X}_i)_{i\in I}\)，它覆盖 \(\mathbf{X}\)，并且对所有 \(i\in I\)，有一个左作用律片段 \(\psi_i\)，它是 \(\mathbf{C}^{r-1}\) 类的，由 \(\mathbf{G}\) 定义在 \(\mathbf{X}_i\) 上，使相应的无穷小作用律为 \(a\mapsto \mathrm{D}_a|\mathbf{X}_i\)。由于 \(\mathbf{X}\) 是旁紧的，可以设覆盖 \((\mathbf{X}_i)_{i\in I}\) 是局部有限的。对所有 \((i,j)\in I\times I\) 及所有 \(x\in \mathbf{X}_i\cap \mathbf{X}_j\)，\(\psi_i\) 和 \(\psi_j\) 在 \((e,x)\) 的某个邻域上重合（命题 11 的推论）。由于 \(\mathbf{X}\) 是正规空间，我们可以应用 § 1, no. 11 的 Proposition 24，从而证明 (i)。

推论 2. 设 X 是 \(\mathbf{C}^r\) 类旁紧流形（\(r \geqslant 2\)），\(\xi\) 是 X 上的 \(\mathbf{C}^{r-1}\) 类向量场。存在 K 在 X 上的作用律片段 \(\psi\)，它是 \(\mathbf{C}^{r-1}\) 类的，使得对所有 \(x \in \mathbf{X}\)，\(\xi(x)\) 是 K 在 0 处的切向量 1 在 \(t \mapsto (t, x)\) 下的像。具有上述性质的两个作用律片段在 \(\{0\} \times \mathbf{X}\) 的某个邻域上重合。

这是推论 1 的特殊情形。

注记。在本号中，左作用律当然可以处处替换为右作用律。

